\chapter{Carrying Specifications: Embedding in Artefacts and Interfaces}\label{ch:article2}

A specification that cannot travel with the artefact it describes cannot become
part of a software lifecycle: it lives in a separate document that drifts out of
date and is unavailable to the very consumers --- test generators, IDEs,
verifiers, and large language models reading an API at runtime --- that would
benefit from it. Having shown in \cref{ch:article1} that inferred contracts earn
their keep downstream, this chapter asks how such a contract can be made to
\emph{travel with} its code, and so bridges the act of \emph{producing}
specifications with the act of \emph{integrating} them into the development
process in \cref{ch:article4}.

\section{The Carriage Problem}
\label{sec:a2-intro}

The inference instrument of \cref{ch:article1} assumes, like much work on
inferred specifications for verification and tooling, that the consumer has
access to the same source tree the inferrer ran over. That assumption is rarely
satisfied in practice. Java libraries are distributed as JAR files; REST
endpoints expose only their HTTP signature; gRPC services expose only their
protobuf descriptors. Source code is the exception, not the norm. If inferred
Java Modeling Language (JML) contracts~\cite{leavens2006jml} are to be consumed
where they are needed, they must be carried by the compiled artefacts and the
interface documents that developers already exchange.

The conceptual question of whether specifications \emph{should} travel alongside
compiled artefacts has a long history in the design-by-contract
tradition~\cite{meyer1992dbc}: a contract is a property of the supplier as seen
by the client, and the client most often holds only the binary. The practical
question, and the one this chapter answers, is whether the standard toolchain
can be made to carry contracts without modification. I propose and evaluate a
uniform mechanism for distributing inferred JML contracts in compiled Java
libraries and in REST service interfaces: a runtime-retained annotation written
into class-file attributes by an ASM-based embedder, plus a parallel
OpenAPI~3.x extension family for service boundaries. The contracts are
consumable by reflection, by any class-file reader, and --- for verification ---
by OpenJML~\cite{cok2014openjml} through a stub-file sidecar.

The study is organised around five questions. First, can inferred JML contracts
be embedded in compiled Java artefacts and reproduced losslessly from the
embedded bytes, with no source-code dependency? Second, what is the
bytecode-size overhead of the embedding, and how does it scale across libraries
of different shapes (utility code, I/O abstractions, generics-heavy
collections)? Third, which format-design choices most reduce that overhead, and
what is the per-choice impact on real inferred specifications versus synthetic
ones? Fourth, is the embedding and reading throughput sufficient to support
per-pull-request continuous-integration workflows on production-scale libraries?
Fifth, does an OpenAPI extension family conform to the OpenAPI~3.x extension
rule, roundtrip losslessly through a reader/writer pair, and survive the
canonical \texttt{swagger-parser} unchanged?

I evaluate the bytecode side empirically on six real-world Java libraries
(Commons Lang, IO, and Math, Guava, jOOL, and Vavr; 39{,}268 methods) and the
full inferrer-to-embedder pipeline on real inferred specifications across the
same six libraries' source trees (22{,}881 specifications). The OpenAPI side is
established by construction against the OpenAPI~3.x \texttt{x-*} extension rule
and empirically by parsing a document carrying every \texttt{x-jml-*} key
through \texttt{swagger-parser}~2.1.x. In summary, the chapter measures
100.0\% lossless roundtrip on every library-regime run; 4.80--6.97\%
bytecode overhead under synthetic specifications and 4.63--10.20\% under real
inferred specifications; a format-optimisation step that reduces overhead by
31\% on real inference and 48--55\% on synthetic specifications; and embedding
throughput of 6{,}000--34{,}000 methods per second, with reading throughput an
order of magnitude higher.

\section{Format Design}
\label{sec:a2-format}

The format is the contract between the embedder and any future consumer. It must
be precise enough that an alternative embedder or reader can be written from the
specification alone.

\subsection{The \texttt{@JmlSpecs} Annotation Type}

A single Java annotation per method carries every clause kind as a separate
\texttt{String[]} member:

\begin{lstlisting}
package com.jml.spec;

@Retention(RetentionPolicy.RUNTIME)
@Target({ElementType.METHOD, ElementType.CONSTRUCTOR,
         ElementType.FIELD, ElementType.TYPE})
public @interface JmlSpecs {
    String[] requires() default {};
    String[] ensures() default {};
    String[] assignable() default {};
    String[] loopInvariant() default {};
    /** Each entry is "ExceptionType|condition". */
    String[] signals() default {};
    String version() default "2";
}
\end{lstlisting}

Within a member, order is the array order; \texttt{requires} clauses are emitted
in source order, which is load-bearing for well-definedness chains (a
precondition asserting an array element's value depends on earlier preconditions
establishing the array's non-nullness and the index's in-bounds). Across
members, the writer fixes a canonical declaration order ---
\texttt{requires}, \texttt{ensures}, \texttt{assignable}, \texttt{loopInvariant},
\texttt{signals}, \texttt{version} --- so two writers that follow this format
produce byte-identical annotation entries for the same logical specification,
even though the JVM annotation format itself does not constrain member order.
Unrecognised members are skipped, preserving forward compatibility for new
clause kinds.

\paragraph{Default-omission convention.}
The \texttt{assignable} member is omitted when its only entry is
\texttt{\textbackslash nothing}; an absent \texttt{assignable} is interpreted as
\texttt{\textbackslash nothing}. This convention exploits a strong empirical
regularity --- 100\% of specifications in the corpus carry
\texttt{assignable \textbackslash nothing}, the inferrer's default for
pure-looking methods --- and removes one array entry per specification. A
specification that does have non-default assignable clauses lists them
explicitly.

\paragraph{Token-encoded clause strings.}
Clause text is encoded before being written to the constant pool. A twelve-entry
dictionary substitutes common JML tokens (\texttt{\textbackslash result},
\texttt{\textbackslash old(}, \texttt{\textbackslash forall},
\texttt{\textbackslash exists}, \texttt{\textbackslash sum},
\texttt{\textbackslash product}, \texttt{\textbackslash num\_of},
\texttt{ ==> }, \texttt{!= null}, \texttt{== null},
\texttt{\textbackslash nothing}, \texttt{\textbackslash everything}) with single
Unicode control bytes (U+0001 through U+0011). Twelve is the empirical sweet
spot: every token in the dictionary represents at least 0.5\% of corpus bytes on
the real-inference corpus, and a learned-dictionary experiment that grows the
dictionary to 28 entries (\cref{subsec:a2-format-opt}) recovers only a further
0.59 percentage points. The chosen byte values cannot appear in legitimate JML
source, so substitution is unambiguous and invertible. UTF-8 encodes each token
byte as a single byte, so each occurrence saves \emph{token-length} $-$ 1 bytes.
The reader decodes on read; consumers that do not know about the dictionary will
see the control characters and can either decode using the published table or
fall back to the sidecar.

\paragraph{Backward and forward compatibility.}
A predecessor format (v1) used Java's repeatable-annotation mechanism: one
\texttt{@JmlSpec} per clause with per-clause \texttt{kind}, \texttt{text},
\texttt{order}, and \texttt{version} members, wrapped in
\texttt{@JmlSpecs(value=\{\dots\})}. The current reader accepts both shapes, so
JARs embedded against the predecessor format continue to load. Only the current
format (v2) is emitted by the writer. The forward-compatibility policy mirrors
the Java language's stance on annotation evolution: unknown annotation members
are skipped by the reader, and token bytes that a future writer adds but a
current reader does not know about pass through as control characters (and
remain decodable by any consumer holding the published dictionary). New clause
kinds can therefore be added without breaking older readers; new dictionary
entries are similarly additive.

\paragraph{Synthetic carriers.}
Lambdas, method references, anonymous inner classes, and \texttt{record}
accessors compile to synthetic methods or classes; the embedder writes
specifications to the synthetic carrier under its bytecode descriptor.

\subsection{Canonical Form}

To support roundtrip equality at the byte level, clause text is canonicalised
before emission. Whitespace runs collapse to single spaces; there are no spaces
around \texttt{.} or \texttt{::}; a single space surrounds binary operators
(\texttt{+}, \texttt{-}, \texttt{*}, \texttt{/}, \texttt{\%}, comparison,
logical); no space appears inside parentheses or brackets; and \texttt{=>} and
\texttt{<=>} are normalised to the long forms \texttt{==>} and \texttt{<==>}.
Quantifier headers retain their bound-variable names; alpha-renaming is deferred
to a future release and is currently a no-op.

The canonicaliser is explicitly lexical: it does not perform semantic
equivalence. Two clauses that differ in operand order ($a > b$ versus $b < a$)
remain distinct strings, and roundtrip equality is preserved on the wire form
rather than on a normalised theory level. Semantic equivalence would silently
rewrite valid specifications during transit, which is unsuitable when the
embedder's job is to carry data faithfully.

\subsection{The OpenAPI Extension}

For service boundaries, a parallel family of OpenAPI~3.x \texttt{x-*} extension
keys carries the same information per operation and per schema:

\begin{lstlisting}[language={}]
paths:
  /orders/{id}:
    get:
      x-jml-requires:
        - "id != null"
        - "id > 0"
      x-jml-ensures:
        - "\\result != null ==> \\result.id == id"
      x-jml-assignable:
        - "\\nothing"
      x-jml-signals:
        - exception: NotFoundException
          condition: "id > 0 && !exists(id)"
\end{lstlisting}

Each extension key is an array of canonical-form strings.
\texttt{x-jml-requires} preserves order; the others are unordered.
\texttt{x-jml-signals} is an array of \texttt{\{exception, condition\}} objects,
with exception types given as JVM internal names so cross-language consumers can
resolve them. A JSON Schema document accompanies the artefact and may validate
any OpenAPI document independently of the inferrer. When the OpenAPI document is
generated by a framework (springdoc-openapi, JAX-RS) and is not editable, a
sidecar JSON file (\texttt{<service>.contract.json}) carries the same content
keyed by route and verb.

\subsection{Sidecar JAR}

For specifications that exceed the per-class bytecode budget --- the JVM
constant-pool cap is 65{,}535 entries per class, though in practice no class
approached this limit in any reported library --- and for signed JARs that
cannot be re-signed, the embedder emits a parallel artefact with the Maven
classifier \texttt{-jml}. The sidecar contains JML stub files
(\texttt{<class>.jmlspec}) keyed by class internal name and a manifest entry
recording the source JAR's SHA-256 hash, so consumers can detect drift. Stub
files are the format OpenJML~\cite{cok2014openjml} reads natively, so the sidecar
is consumable by the verifier without further transformation.

\section{Implementation}
\label{sec:a2-impl}

The embedder is implemented in Java~21, using ASM~9.7 for class-file rewriting.
The module \texttt{com.jml:jml-embedder} is approximately 1{,}800 source lines of
code, of which roughly 1{,}000 are in the three core classes described below; the
remainder cover the codec, the OpenAPI extension reader/writer, an automated
dictionary learner used in \cref{subsec:a2-format-opt}, and small data records.
The module depends only on ASM, Gson, and SLF4J.

The data flow is a write path and a symmetric read path. The predecessor
inferrer of \cref{ch:article1} produces a \texttt{MethodSpecification}; a bridge
class, \texttt{InferrerSpecConverter}, translates it to the embedder's
\texttt{MethodSpec} record byte-for-byte, preserving clause-list order;
\texttt{AsmJmlSpecWriter} (with the codec and canonicaliser) emits an annotated
JAR; and \texttt{AsmJmlSpecReader} recovers the original \texttt{MethodSpec} map.

\subsection{The Writer}

The writer accepts either a single class file (as \texttt{byte[]}) or a JAR (as
a \texttt{Path}). For a JAR, it walks each \texttt{.class} entry, applies a
\texttt{ClassReader} $\to$ \texttt{ClassVisitor} $\to$ \texttt{ClassWriter}
pipeline that intercepts \texttt{visitMethod} calls, and emits a single
\texttt{@JmlSpecs} annotation per method with each clause kind as a
\texttt{String[]} member. Order within each kind is preserved by array order.
Clause strings are encoded via the dictionary codec before being written, and
the default \texttt{assignable \textbackslash nothing} member is omitted. Methods
without a corresponding \texttt{MethodSpec} are passed through unchanged. For
specifications that exceed the per-class bytecode budget, or for JARs that must
remain unmodified, the writer also produces the sidecar JAR; its stub files are
generated by a small printer that decodes JVM method descriptors back into
source-level types.

\subsection{The Reader}

The reader is symmetric: it walks each \texttt{.class} entry with a
\texttt{ClassReader} configured to skip code, debug, and frames (it needs only
the annotations) and accumulates a \texttt{Map<MethodKey, MethodSpec>}. It
accepts both the current v2 shape (a single \texttt{@JmlSpecs} with array members
per clause kind) and the legacy v1 shape (repeated \texttt{@JmlSpec} annotations
with per-clause metadata). For v2, clause strings are decoded via the codec, and
a missing \texttt{assignable} member is synthesised as
\texttt{[\textbackslash nothing]}.

\subsection{The Canonicaliser and Pipeline Integration}

The canonicaliser implements the rules of \cref{sec:a2-format}. It tokenises the
clause text with a small lexer (identifiers, numeric literals, multi-character
operators, JML backslash-keywords) and re-emits the tokens with canonical
spacing. The implementation is 260 source lines, validated by 16 test cases
covering whitespace collapsing, operator spelling, parenthesisation, and
idempotency. With the converter in place, the end-to-end pipeline is one method
call:

\begin{lstlisting}
MethodSpecification inferred = new MethodSpecificationInferrer()
        .inferSpecification(method);
MethodSpec embeddable = InferrerSpecConverter.toEmbeddable(inferred);
new AsmJmlSpecWriter().embedJar(input, output,
        Map.of(methodKey, embeddable));
\end{lstlisting}

The roundtrip property
\texttt{extract(embed(infer(source))) == infer(source)} is exercised by a JUnit
suite at small scale and at corpus scale by the validation harness of the next
section.

\section{Empirical Validation}
\label{sec:a2-empirical}

The embedder is evaluated on real-world Java libraries. The harness applies the
embedder to each library's JAR, measures four quantities --- lossless-roundtrip
rate, bytecode-size overhead, embedding throughput, and reading throughput ---
and does so under two regimes: synthetic specifications that isolate the
embedding mechanism, and real inferred specifications that exercise the full
pipeline.

\subsection{Subjects and Method}

Six libraries are reported: Apache Commons Lang~3.14.0, Apache Commons
IO~2.13.0, Apache Commons Math~3.6.1, Guava 33.3.0-jre, jOOL~0.9.14, and
Vavr~0.10.4. Commons Lang is the subject used by the inference study of
\cref{ch:article1}; Commons IO contributes file/stream utility classes with
significant abstract-method counts; Commons Math adds numerical-utility code;
Guava is a larger contrast with more complex generics and a multi-release JAR;
and jOOL and Vavr contribute functional-programming idioms (heavy use of generic
interfaces, default methods, and lambdas). Together the six libraries cover
39{,}268 methods.

For the synthetic regime, the harness generates one specification per method
(one \texttt{requires} clause per non-primitive parameter, one \texttt{ensures}
clause for reference returns, and \texttt{assignable \textbackslash nothing}),
embeds the specification map into a copy of the JAR, reads it back, and matches
specifications against the originals on \texttt{(internalName, methodName,
descriptor)} keys. Both original and embedded JAR sizes are on-disk file sizes
measured after the DEFLATE compression that the \texttt{.jar} container applies,
so the reported overhead already accounts for any constant-pool string
redundancy the compressor can exploit. A separate negative test loads an
embedded class via a \texttt{URLClassLoader} that does not see the annotation
type and confirms that class load succeeds: reflective annotation access is
allowed to raise \texttt{TypeNotPresentException}, but class load itself must not
fail.

\subsection{Results}

\Cref{tab:a2-validation} reports the headline measurements. Method counts,
lossless-roundtrip rate, and bytecode-size overhead are deterministic and
reported once; throughput figures vary across runs (JIT warm-up, garbage
collection, OS scheduling) and are reported as medians across 39 independent
harness invocations per library.

\begin{table}[ht]
\centering
\caption{Embedder validation on real-world libraries (medians across 39 runs for
throughput; lossless-roundtrip rate and overhead are deterministic). The lower
block reports the full inferrer-to-embedder pipeline on real inferred
specifications; the synthetic specifications in the upper block isolate the
embedding mechanism. On matched methods, equality is exact (zero clause-level
mismatches).}
\label{tab:a2-validation}
\begin{tabular}{lrrrrr}
\toprule
\textbf{Library} & \textbf{Methods} & \textbf{Lossless} & \textbf{Byte ovhd.} & \textbf{Embed (m/s)} & \textbf{Read (m/s)} \\
\midrule
\multicolumn{6}{l}{\textit{Synthetic specs (isolating the embedding mechanism)}} \\
Apache Commons Lang 3.14.0  & 4{,}029  & 100.0\% & +5.85\% & 16{,}225 & 221{,}617 \\
Apache Commons IO 2.13.0    & 2{,}677  & 100.0\% & +6.97\% & 21{,}976 & 327{,}618 \\
Apache Commons Math 3.6.1   & 9{,}191  & 100.0\% & +5.19\% & 20{,}485 & 260{,}593 \\
Guava 33.3.0-jre            & 13{,}840 & 100.0\% & +5.27\% & 12{,}618 & 127{,}924 \\
jOOL 0.9.14                 & 3{,}869  & 100.0\% & +4.80\% & 30{,}218 & 226{,}339 \\
Vavr 0.10.4                 & 5{,}662  & 100.0\% & +5.56\% & 34{,}138 & 333{,}994 \\
\midrule
\multicolumn{6}{l}{\textit{Real inferred specs (full inferrer-to-embedder pipeline)}} \\
Apache Commons Lang 3.14.0  & 2{,}332  & 100.0\% & +10.07\% & 7{,}680  & 78{,}801 \\
Apache Commons IO 2.13.0    & 1{,}680  & 100.0\% & +10.20\% & 10{,}702 & 101{,}218 \\
Apache Commons Math 3.6.1   & 6{,}001  & 100.0\% & +10.04\% & 7{,}595  & 111{,}642 \\
Guava 33.3.0-jre            & 6{,}320  & 100.0\% & +4.63\%  & 5{,}951  & 68{,}281 \\
jOOL 0.9.14                 & 2{,}648  & 100.0\% & +5.81\%  & 13{,}668 & 126{,}159 \\
Vavr 0.10.4                 & 3{,}900  & 100.0\% & +6.39\%  & 14{,}398 & 133{,}762 \\
JML-Inferrer self-test      & 73       & 100.0\% & ---      & ---      & --- \\
\bottomrule
\end{tabular}
\end{table}

\paragraph{Lossless roundtrip.}
All six libraries achieve a 100.0\% lossless-roundtrip rate. An earlier prototype
reported 96.9\% / 97.6\% on Commons Lang and Guava owing to a missed callback in
the writer's \texttt{MethodVisitor} lifecycle: lazy emission was triggered by
\texttt{visitCode}, \texttt{visitParameter}, \texttt{visitAnnotationDefault}, or
\texttt{visitAnnotation}, but none of these fires for abstract or interface
methods, which only invoke \texttt{visitEnd}. Adding \texttt{visitEnd} to the
trigger set closed the gap, so every method on every interface in all six
libraries now receives the embedded annotation. Commons IO has a noticeably
higher proportion of interfaces, which is why its byte-overhead percentage runs
higher than the others under both formats.

\paragraph{Byte overhead.}
The synthetic-spec overhead sits in the 4.80--6.97\% range. The differences are
explained by structural shape rather than total size: Commons IO has the highest
proportion of interfaces (each of which receives a clause set per method),
pushing its overhead above the others; jOOL has the lowest, consistent with its
smaller per-class method counts and concentration of simple functional-interface
declarations. The absolute overhead per library is well below typical Maven
distribution sizes (for example, 38~KB for Commons Lang, 162~KB for Guava). The
synthetic figures are not, however, an upper bound on real overhead. The
real-inference run on the same Commons Lang JAR measures 10.07\% --- 4.22
percentage points above the synthetic 5.85\% --- because real inferred
specifications are richer than the synthetic worst case: 90.4\% have at least one
\texttt{ensures} clause, 11.5\% carry one or more loop invariants (the synthetic
harness emits none), 14.8\% include branch-conditional implications, and the
average clause string is longer because the inferrer emits parameter bounds,
\texttt{\textbackslash result} relations, and \texttt{\textbackslash old(\dots)}
references that the synthetic policy never produces. A deployer planning a
storage budget should provision for the real-inference figure (around 10\% on a
typical utility library), not the synthetic-uniform figure.

\paragraph{Throughput.}
Synthetic-spec median embedding throughputs span 12{,}618--34{,}138 methods per
second (Guava slowest, Vavr fastest); reading throughputs span
127{,}924--333{,}994 methods per second. Real inference is slower because each
specification carries more clauses: 5{,}951--14{,}398 m/s embedding,
68{,}281--133{,}762 m/s reading. At the synthetic median for Guava, embedding the
full 13{,}840-method JAR takes roughly 1.1~seconds; under real inference,
2.3~seconds. Reading is consistently an order of magnitude faster because the
reader skips full class rewriting. Both regimes comfortably support
per-pull-request continuous-integration workflows.

\paragraph{Negative test.}
The class \texttt{org.apache.commons.lang3.StringUtils} loaded successfully via a
\texttt{URLClassLoader} that does not see the annotation type. Reflective
annotation access either returns the resolvable annotations (none in this
loader's view) or raises \texttt{TypeNotPresentException}, both documented
behaviours; class load is unaffected. An existing JAR can therefore be enriched
in place with no risk to consumers unaware of the annotation, and with no
build-toolchain changes.

\subsection{Real Inferred Specifications Across Six Libraries}
\label{subsec:a2-real-inference}

The synthetic-spec validation isolates the embedding mechanism from the
inferrer's behaviour, but does not exercise the inferrer-to-embedder pipeline on
real code. A second harness closes that gap: it runs the inferrer over every
method in the source trees of all six libraries, resolves the source signatures
against each binary JAR's bytecode descriptors by name and parameter count,
embeds the resulting specifications, reads them back, and checks clause-list
equality.

Across the six libraries the harness processes 2{,}339 source files containing
32{,}459 source method declarations. The source-to-bytecode key resolution
places 22{,}881 specifications into distinct keys (1{,}680 Commons IO + 2{,}332
Commons Lang + 6{,}001 Commons Math + 6{,}320 Guava + 2{,}648 jOOL + 3{,}900
Vavr), and every one of them roundtrips with byte-for-byte equality on its clause
list: matched equals total, zero missing, zero mismatches in every library. The
real-inference byte-overhead figures sit at 4.63--10.20\%. The Apache libraries
(Lang 10.07\%, IO 10.20\%, Math 10.04\%) cluster at the upper end because real
specifications are richer than the synthetic worst case; the
functional-programming libraries (jOOL 5.81\%, Vavr 6.39\%) and Guava (4.63\%)
sit lower because the same per-spec byte cost amortises against larger, denser
classes.

A residual of source-method-to-bytecode-descriptor lookups fails on overload
collisions, where two source overloads of the same name and arity map to the
same bytecode descriptor under the harness's first-match-wins choice. Roughly one
quarter of inferred-with-clauses methods are lost this way on Commons Lang, fewer
on the others. The 100\% lossless rate is on every specification the harness
successfully embeds; a full type-resolving harness would recover the residual
without changing the embedder's behaviour.

\paragraph{Spec-shape coverage.}
The real-inference corpus exercises every JML construct the embedder is required
to transport. \Cref{tab:a2-shapes} reports the distribution.

\begin{table}[ht]
\centering
\caption{Spec-shape coverage aggregated across all six libraries' 22{,}881
real-inference specifications. Every category roundtrips losslessly on every
library.}
\label{tab:a2-shapes}
\begin{tabular}{lrr}
\toprule
\textbf{Construct} & \textbf{Specs} & \textbf{\% of corpus} \\
\midrule
\texttt{requires} clauses               & 9{,}734  & 42.5\% \\
\texttt{ensures} clauses                & 20{,}601 & 90.0\% \\
\texttt{assignable} clauses             & 22{,}881 & 100.0\% \\
\texttt{loop\_invariant} clauses        & 2{,}458  & 10.7\% \\
Quantifiers (\texttt{\textbackslash forall}, \texttt{\textbackslash exists}, \texttt{\textbackslash sum}) & 312 & 1.4\% \\
\texttt{\textbackslash old} references  & 411    & 1.8\% \\
\texttt{\textbackslash result} references & 15{,}809 & 69.1\% \\
Branch-conditional (\texttt{==>})       & 1{,}958  & 8.6\% \\
\bottomrule
\end{tabular}
\end{table}

Per-library proportions vary as expected: Commons Math has the highest
loop-invariant share (20.9\%) and quantifier share (3.8\%); jOOL and Vavr have
the highest \texttt{ensures}-clause shares (97.7\%, 99.3\%); Commons IO has the
lowest \texttt{ensures} share (76.3\%, dominated by void stream utilities).
Despite this variation, every category roundtrips losslessly across all six
libraries, so the embedder transports not only simple null-check preconditions
but also nested quantifiers, \texttt{\textbackslash old} references, conditional
implications, and loop invariants. A 73-method self-test on the inferrer's own
source code supplements this with an end-to-end check on a smaller and
stylistically different corpus, and also achieves a 100\% lossless-roundtrip
rate.

\subsection{Format-Optimisation Impact}
\label{subsec:a2-format-opt}

The earlier v1 format used Java's repeatable-annotation mechanism --- each clause
produced a separate \texttt{@JmlSpec} annotation with \texttt{kind},
\texttt{text}, \texttt{order}, and \texttt{version} members, wrapped in a
\texttt{@JmlSpecs(value=\{\dots\})} container. This maximally verbose shape
produced 11.76--13.49\% overhead on the synthetic corpora and 14.70\% on the
real-inference Commons Lang run. The current v2 format applies three changes: a
single annotation per method (eliminating per-clause annotation framing),
default-omission for \texttt{assignable \textbackslash nothing}, and the twelve-token
dictionary. \Cref{tab:a2-opt} reports the impact; the lossless-roundtrip property
is preserved at 100\% under both formats.

\begin{table}[ht]
\centering
\caption{Format-optimisation impact on bytecode-size overhead. The
lossless-roundtrip rate is 100\% under both formats.}
\label{tab:a2-opt}
\setlength{\tabcolsep}{3pt}
\begin{tabular}{lrrr}
\toprule
\textbf{Corpus} & \textbf{v1 ovhd.} & \textbf{v2 ovhd.} & \textbf{Reduction} \\
\midrule
Synthetic --- Commons Lang      & 11.95\% & 5.85\%  & $-$51\% \\
Synthetic --- Commons IO        & 13.49\% & 6.97\%  & $-$48\% \\
Synthetic --- Guava             & 11.76\% & 5.27\%  & $-$55\% \\
\midrule
Real inference --- Commons Lang & 14.70\% & 10.07\% & $-$31\% \\
\bottomrule
\end{tabular}
\end{table}

The synthetic-spec reduction is larger than the real-inference reduction because
synthetic specifications are dominated by exactly the patterns the optimisation
targets: every specification carries the default-omittable
\texttt{assignable \textbackslash nothing}, every \texttt{requires} clause is of
the form \texttt{p != null} (a dictionary token), and every \texttt{ensures}
clause references \texttt{\textbackslash result} (also a token). The
real-inference corpus contains a wider variety of clause shapes; many strings
still benefit from the dictionary, but a long tail of arithmetic expressions,
range comparisons, and quantified subexpressions falls outside the current twelve
tokens. The real-inference reduction is therefore a more realistic estimate of
what a production deployment should expect.

The choice of which JML tokens to encode and which clause to declare a default is
not arbitrary: each rests on the shape distribution measured on the
real-inference corpus (\cref{tab:a2-shapes}). A format designed without the
measurement would either include every JML keyword in the dictionary (paying
constant-pool cost for tokens that never fire) or omit tokens that occur
frequently in real specifications. The synthetic-versus-real divergence
(48--55\% reduction versus 31\%) is itself empirical evidence that the saving
depends on which clause shapes the consumer's inferrer produces.

A configuration sweep on the 2{,}332-spec Commons Lang real-inference corpus
confirms the contribution of each convention, with every mode roundtripping
losslessly. Disabling the codec (\texttt{WriterConfig.NONE}) raises overhead to
11.26\%, so the default-omission and token-dictionary conventions together
account for 1.19 percentage points of saving over an always-explicit baseline.
An automated dictionary learner that extends the dictionary from 12 to 28 tokens
gains a further 0.59 percentage points (down to 9.48\%); the improvement is real
but modest, establishing that on the tested corpora a hand-picked-good dictionary
is close to learned-best. A compressed v3 mode that DEFLATEs each method's clause
set into a single \texttt{byte[]} annotation member is a clear negative result at
30.25\% overhead: DEFLATE on 50--100-byte payloads cannot exploit much
redundancy, and the per-spec compressed payloads are unique enough that the
constant-pool deduplication v2 enjoys becomes unavailable. The v3 mode is
retained as an opt-in for completeness but is not recommended as the default.

\section{Discussion}
\label{sec:a2-discussion}

\paragraph{Why an annotation, not a custom attribute?}
Custom class-file attributes are a tempting alternative: they avoid the
constant-pool pollution of large string literals and are the more general
mechanism. Three reasons argue against them in practice. First, tool visibility:
standard reflection ignores custom attributes, IDEs do not display them, and
shrinkers such as ProGuard and R8 strip unrecognised attributes by default, so
each downstream consumer would need a custom reader. Second, repeatable
annotations make per-clause emission straightforward, removing the historical
reason to prefer a custom attribute. Third, an annotation carrier with clause
kinds as named members evolves more compatibly than a custom attribute's binary
layout. The annotation is therefore the primary form; the custom-attribute route
and the sidecar JAR remain available where the annotation is unsuitable (signed
JARs, aggressive obfuscation).

\paragraph{Generality.}
The format is independent of the inferrer that produces the specifications. Any
source of JML clauses --- the inferrer of \cref{ch:article1}, a dynamic invariant
detector, hand-written contracts, or contracts produced by a future
LLM-assisted refinement step --- can be embedded by the same mechanism. The
format also generalises to other JVM languages (Kotlin, Scala) at the cost of
those languages supplying their own translators between source-level annotations
and the bytecode representation. For non-JVM languages, the OpenAPI extension
carries equivalent information at the service boundary: the bytecode side is
JVM-specific, the service-boundary side is language-agnostic.

\paragraph{Limitations.}
The embedder is consumer-tolerant, but the consumer has work to do. A reader that
wants to act on the embedded clauses must parse the canonical-form JML strings,
substitute parameter names, and decide what to do with the result. This chapter
ships the writer, the reader, and the canonicaliser, but not a verifier or
runtime checker --- those are downstream concerns that benefit from the carriage
mechanism without being implemented by it. One substantive limitation is that
roundtrip equality is on the canonical-form string, not on a normalised semantic
theory, so equivalent specifications written in different forms appear distinct
on the wire. This separation of concerns is deliberate; semantic normalisation is
a separable layer.

\section{Threats to Validity}
\label{sec:a2-threats}

\paragraph{Internal validity.}
The 100\% lossless-roundtrip rate depends on the writer covering every method
shape that ASM emits; the prototype's 96.9\% / 97.6\% gap, traced to the missing
\texttt{visitEnd} trigger for abstract and interface methods, shows the fragility
of this dependence. The expanded six-library corpus exercises
functional-programming idioms and numerical-utility code in addition to the
original three subjects, and the embedder reaches 100\% on every category. The
same risk pattern applies to method shapes specific to future JVM features
(abstract record components, hidden classes), which would need their own triggers
if they exhibit different visit-call orderings; the mitigation is a corpus-scale
validation per JVM release. The byte-overhead figures are deterministic given the
JAR and the spec map and reproduce byte-for-byte on replication. Throughput
figures vary across runs and were collected on a single platform (OpenJDK~21 with
default G1GC on a Windows~11 workstation, no concurrent load) across 39
invocations per library, with raw per-run numbers retained for re-analysis.

\paragraph{External validity.}
The six reported libraries are mature, well-structured, and shipped through Maven
Central; generalisation to less-tidy code (older artefacts, vendored
dependencies, hand-tuned shrinker-stripped JARs) needs separate validation. The
negative test confirms the consumer-side contract at JDK~21; older JDKs and
ahead-of-time compilation paths are deferred. The real-inference validation
covers all six libraries (22{,}881 specifications in total), settling the
cross-corpus generalisation question for the lossless-roundtrip claim across the
reported corpus. The OpenAPI extension is validated empirically against the
canonical OpenAPI~3.x parser, \texttt{swagger-parser} 2.1.x: a JUnit suite parses
a document carrying every \texttt{x-jml-*} key and asserts that the in-memory
model exposes each extension on the operation object with key names and value
lists preserved, including the nested \texttt{\{exception, condition\}} structure
of the \texttt{signals} clause. The OpenAPI~3.x specification requires
\texttt{x-*} preservation, and the test confirms the reference implementation
honours that requirement for this extension family. End-to-end validation against
\texttt{openapi-generator} and a real microservice's emitted document is planned
follow-up.

\paragraph{Construct validity.}
The lossless-roundtrip metric is set-equality on \texttt{(internalName,
methodName, descriptor)} keys with byte-for-byte clause-list equality on each
matched key. Two methods with identical signatures but different bytecode bodies
appear identical to the harness; this is the right construct for an embedder that
carries metadata rather than bytecode, but a reader needing to attach annotations
to a specific compilation would require additional disambiguation. The
real-inference harness's name-plus-arity resolution loses information on overload
pairs of the same arity, costing 917 of 3{,}249 inferred-with-clauses methods on
Commons Lang; this is a measurement limitation rather than an embedder defect, as
the embedder roundtrips 100\% of every specification it receives.

\paragraph{Conclusion validity.}
The empirical claims are point measurements on six libraries under both regimes;
no statistical significance test is meaningful, since the byte-overhead figures
are deterministic and the throughput figures are reported as medians across 39
runs to control for run-to-run variability. The conclusions hold for the JARs and
specification sets reported, and the relevant generalisation question is breadth
(more libraries, more JDK versions) rather than statistical inference.

\section{Conclusion}
\label{sec:a2-conclusion}

Automatically inferred JML specifications can be embedded in compiled Java
artefacts and REST API descriptions without loss: a runtime-retained annotation
\texttt{@JmlSpecs} written by an ASM-based embedder, plus a parallel OpenAPI~3.x
extension family. Roundtrip is lossless on six libraries totalling 39{,}268
methods under synthetic specifications and on 22{,}881 real inferred
specifications across the same six libraries' source trees; overhead is
4.63--10.20\% on real inferred specifications after the format-optimisation step,
a 31\% reduction over the v1 per-clause format on the measured corpus; the
OpenAPI extension is preserved unchanged by \texttt{swagger-parser}~2.1.x; and
embedded classes load correctly in JVMs without the annotation type.

This settles the carriage question (RQ3) for the thesis: the path from source
code to consumer-accessible contract is complete and toolchain-transparent, so a
specification produced by the inference of \cref{ch:article1} can now travel with
the artefact a developer actually deploys. With a vehicle for carriage in place,
the thesis turns in \cref{ch:article4} to integrating these contracts into the
development process, where the same carrier transports the compositional
specifications that cross method and library boundaries.
